% !PRINTED 447
\weberpage{447}
From this it now follows again, by (7), that the congruence % [sic] The German print cites (7); the inference rests on the immediately preceding congruences (19)--(20).
\weq{21}{\xi^{p^\varepsilon} \equiv \xi \quad (\mathrm{mod}\ \mathfrak P)}
is satisfied for every integer $\xi$ in $\mathfrak K$, and hence the theorem:

\gesperrt{A prime $p$ splits in the field $\mathfrak K$ into prime ideals all different from one another, whose degree is equal to the index of $p$.}

We wish to prove further that among the different
prime ideals $\mathfrak P$ dividing $p$ there is always one for which,
of the two congruences following from (16), the congruence
(17) holds.

To this end we remark: if we transform all the numbers of the
field $\mathfrak K$ into the conjugate imaginary ones, interchanging
$\sqrt{D}$ with $-\sqrt{D}$ and every class invariant $(k)$ with the opposite one $(k^{-1})$, then all the numbers of an
ideal $\mathfrak A$ pass over into the numbers of a conjugate imaginary ideal $\mathfrak A'$,
which may also be identical with $\mathfrak A$, and the norm of the
ideal $\mathfrak A$ is equal to the norm of the ideal $\mathfrak A'$. If $\mathfrak A$ is a prime ideal,
then $\mathfrak A'$ is a prime ideal too.

If now the congruence (17) is not satisfied, then by (16)
\[
(lk)^p \equiv (k) \quad (\mathrm{mod}\ \mathfrak P)
\]
must hold, and this congruence too remains in force when $k$ is
replaced by another class $k'$. If, then, we put $l^{-1}k^{-1}$ in
place of $k$, it follows that:
\[
(k^{-1})^p \equiv (l^{-1}k^{-1}) \quad (\mathrm{mod}\ \mathfrak P).
\]

Since therefore $(k^{-1})^p - (l^{-1}k^{-1})$ is a number in $\mathfrak P$, $(k)^p - (lk)$
is contained in the ideal $\mathfrak P'$ conjugate to $\mathfrak P$, and we have
\[
(k)^p \equiv (lk) \quad (\mathrm{mod}\ \mathfrak P'),
\]
which arises from (17) by the interchange of $\mathfrak P$ with $\mathfrak P'$.

7. From these theorems the irreducibility of the class equation, even after adjunction of $\sqrt{D}$, is a simple consequence.

If $k, k'$ are any two classes of the discriminant $D$, then
by § 118 we can compose $k'$ in the following manner:
\[
k' = k l l' l'' \ldots,
\]
so that the classes $l, l', l'', \ldots$ represent the primes $p, p', p'', \ldots$.
By what has just been proved, we can therefore choose
the prime divisors $\mathfrak P, \mathfrak P', \mathfrak P''$ of these primes in such a way that
%CONT>>
